% !PRINTED 680
\weberpage{680}
%<<CONT

Es seien $\alpha_0, \beta_0$ die endlichen Werte der Funktionen $\alpha, \beta$ in einem Punkt $\mathfrak{P}$. Dann ist $F(\alpha_0, \beta_0) = 0$, und wir können die Gleichung (1) so darstellen:
\weq{3}{
\begin{aligned}
0 = {} & (\alpha - \alpha_0)F'(\alpha_0) + (\beta - \beta_0)F'(\beta_0) \\
& + \tfrac{1}{2}\left[(\alpha-\alpha_0)^2 F''(\alpha_0,\alpha_0) + 2(\alpha-\alpha_0)(\beta-\beta_0)F''(\alpha_0,\beta_0) \right. \\
& \left. \quad + (\beta-\beta_0)^2 F''(\beta_0,\beta_0)\right] + \cdots
\end{aligned}
}
Die beiden zueinander reziproken Funktionen
\[
\frac{\alpha - \alpha_0}{\beta - \beta_0}, \quad \frac{\beta - \beta_0}{\alpha - \alpha_0}
\]
können nicht beide in dem Punkt $\mathfrak{P}$ unendlich sein. Es sei also:
\[
\left(\frac{\alpha - \alpha_0}{\beta - \beta_0}\right)_0
\]
der endliche Wert des ersten dieser Quotienten. Dann ergibt sich aus (3), wenn $F'(\alpha_0)$ nicht Null ist:
\weq{4}{\left(\frac{\alpha - \alpha_0}{\beta - \beta_0}\right)_0 = -\frac{F'(\beta_0)}{F'(\alpha_0)} = -\left(\frac{F'(\beta)}{F'(\alpha)}\right)_0 = \left(\frac{d\alpha}{d\beta}\right)_0.}
Diese Gleichung besteht für alle Punkte $\mathfrak{P}$, mit etwaiger Ausnahme einer endlichen Anzahl\footnote{Unter den auszunehmenden Punkten sind jedenfalls die enthalten, in denen $\alpha_0$ oder $\beta_0$ unendlich wird. Denn für diese Punkte sind die Funktionen $\alpha - \alpha_0$, $\beta - \beta_0$ gar nicht erklärt.}.

1. \gesperrt{Eine Funktion $\eta$ in $\Omega$, die in allen Punkten, mit Ausnahme einer endlichen Anzahl, der Bedingung genügt:}
\weq{5}{\eta_0 = \left(\frac{\alpha-\alpha_0}{\beta-\beta_0}\right)_0}
\gesperrt{ist mit $\left(\dfrac{d\alpha}{d\beta}\right)$ identisch.}

Denn aus (4) und (5) folgt, daß die Funktion
\[
\left(\frac{d\alpha}{d\beta}\right) - \eta
\]
unendlich viele Nullpunkte hat und daher identisch verschwinden muß.

Die Formel (5) kann also gleichfalls als Definition des Differentialquotienten dienen.

Daraus ergeben sich einfach die Hauptsätze über die Differentialquotienten:
%CONT>>
